% Subseção da F-RES-2 (ranking de criticidade, RF-RES-04). Dono da sessão: Henrique.
% Fontes: docs/plano-de-desenvolvimento/features/feature-F-RES-2.md (Como funciona e
% Resultados). Números: benchmarks/resultados/falhas-extensoes-2026-11-20.json (comando no JSON).
\subsection{Ranking de criticidade}\label{sec:F-RES-2}

Simular a falha de todas as $V$ paradas custaria $V$ simulações; o ranking só simula
candidatos: as $N$ paradas de maior intermediação (Seção~\ref{sec:F-MET}), que concentram
caminhos mínimos, e as linhas (todas ou uma amostra sorteada). Cada candidato falha sozinho
sobre a rede intacta, com as mesmas $s$ origens sorteadas e o mesmo denominador, e os
candidatos são ordenados pela perda de eficiência. A linha de base é medida uma vez: o
custo é $O((c + 1)\,s\,(V + E)\log V)$ para $c$ candidatos, mais o Brandes amostrado.
Testes com gabarito à mão confirmam que articulações e pontes lideram e que cada perda é
igual à de uma simulação isolada. No feed real ($V = 9{.}660$, $E = 122{.}799$), com as
10 paradas mais centrais, 20 linhas sorteadas e 30 origens, o ranking dos 30 candidatos
levou 38,9~s (mais 4,0~s do Brandes com 100 origens). Lideram três paradas vizinhas (Rua
Deputado Salim Nacur, 173, e Rua das Canárias, 1105 e 729), com perdas de 1,11\%, 1,07\% e
0,89\%; a linha mais crítica, a 9550, fica em 5.º lugar (0,40\%). Nenhum candidato tirou
paradas da rede principal, e a parada de maior intermediação causa só 0,14\%: centralidade
e criticidade não coincidem. \textbf{Limitação:} a intermediação é um filtro, não
uma garantia; uma parada pouco central que desconecta um ramal pode ficar fora do ranking.
